\begin{center}
    \begin{tabular}{ m{2.5cm} m{10cm} m{2.5cm} }
        \includegraphics[width=2.5cm]{images/insat.jpg} &
        \centering
        \begin{minipage}[c]{\linewidth}
            \centering
            {\small \textbf{REPUBLIC OF TUNISIA}}\\[0.15cm]
            {\small Ministry of Higher Education and Scientific Research}\\[0.15cm]
            {\small University of Carthage}\\[0.15cm]
            {\small National Institute of Applied Sciences and Technology}
        \end{minipage}
        &
        \includegraphics[width=2.5cm]{images/unicar.png} \\
    \end{tabular}
\end{center}


\noindent\rule{\textwidth}{0.4pt}
\vspace{0.5cm}
% ------------------- English Abstract -------------------

{\textbf{Design and Implementation of a Comprehensive \acrshort{oee} Framework for an Automated Transport System}}\\[-0.5cm]
\noindent\rule{\textwidth}{0.4pt}
{\textbf{Abstract}}


\small
This graduation project presents an Overall Equipment Effectiveness (\acrshort{oee}) framework adapted to autonomous yard-logistics operations. The approach redefines Availability, Performance, and Quality for variable transport missions and integrates Operational Design Domain (\acrshort{odd}) constraints. A canonical data model and ingestion pipeline were developed to harmonise telemetry, manual logs, and weather data. The framework was implemented at transfer level and extended with mission-phase diagnostics. Industrial data demonstrated the calculation chain, while controlled synthetic scenarios validated the consistency and sensitivity of the implemented metrics.

\textbf{Keywords:} \acrshort{oee}, autonomous vehicle, yard logistics, \acrshort{odd}, canonical data model, mission-phase diagnostics


\vspace{0.5cm}

% ------------------- French Abstract -------------------


{\textbf{Conception et implémentation d’un cadre complet de TRS pour un système de transport automatisé}}\\ [-1cm]

\noindent\rule{\textwidth}{0.4pt}

{\textbf{Résumé}}


\begin{french}
\small
Ce projet présente un cadre de Taux de Rendement Synthétique (TRS) adapté aux opérations de transport autonome en environnement logistique. L’approche adapte la Disponibilité, la Performance et la Qualité aux missions variables et intègre les contraintes \acrshort{odd}. Un modèle de données canonique et une chaîne d’intégration ont été développés pour harmoniser les données télémétriques, les journaux manuels et les données météorologiques. Le cadre est appliqué au niveau des transferts et complété par une analyse des phases de mission.

\textbf{Mots-clés :} TRS, véhicule autonome, logistique de parc, \acrshort{odd}, modèle de données canonique, phases de mission
\end{french}


\vspace{0.5cm}

% ------------------- Arabic Abstract (native XeTeX RTL) -------------------
\artext{{\Large\bfseries تصميم وتنفيذ إطار شامل للفعالية الإجمالية للمعدات لنظام نقل آلي}}
\vspace{-0.3cm}
\noindent\rule{\textwidth}{0.4pt}

\artext{{\bfseries ملخص}}

\small
\artext{يقدّم هذا المشروع إطارًا للفعالية الإجمالية للمعدات (\acrshort{oee}) مكيّفًا مع عمليات النقل الذاتي في الساحات اللوجستية. يعتمد الإطار على تكييف مفاهيم الإتاحة والأداء والجودة مع طبيعة مهام النقل المتغيّرة، مع دمج قيود مجال التشغيل التصميمي (\acrshort{odd}). كما تم تطوير نموذج بيانات موحّد ومسار لمعالجة ودمج بيانات القياس الآلي والسجلات اليدوية والبيانات الجوية. طُبّق الإطار على مستوى عمليات النقل، مع إضافة طبقة تشخيصية على مستوى مراحل المهمة لتحديد مواقع الخسائر بشكل أدق. واستُخدمت البيانات الصناعية لإظهار سلسلة الحساب، بينما استُخدمت سيناريوهات اصطناعية مضبوطة للتحقق من اتساق المؤشرات وحساسيتها.}

\artext{{\bfseries الكلمات المفتاحية:} الفعالية الإجمالية للمعدات، المركبات الذاتية، لوجستيات الساحات، \acrshort{odd}، نموذج البيانات، مراقبة الأداء، مراحل المهمة}

